% ================================================================
% SECTION I: INTRODUCTION
% ================================================================

A blockchain can scale only if every participant can be confident that block data has been published and can be retrieved.
This property, \emph{data availability} (DA), is a security requirement for rollups, parachains and sharded designs.
Without it, validators cannot verify state transitions, light clients cannot detect fraud, and the system's trust model reduces to that of a permissioned committee~\cite{Saif_Survey_DA_Rollups}.

Several DA designs have reached production in the past three years.
Ethereum introduced blob-carrying transactions (EIP-4844) and, with the Fusaka upgrade in December 2025, activated PeerDAS for cell-level sampling~\cite{EIP-4844, EIP-7594}.
Celestia launched the first dedicated DA network with two-dimensional erasure coding and namespace-filtered sampling~\cite{AlBassam_LazyLedger}.
Polkadot's ELVES protocol distributes erasure-coded parachain data across its validator set with approval-gated finality~\cite{Jeff_Burdges_ELVES}.
Espresso's Tiramisu introduces a three-layer Verifiable Information Dispersal (VID) architecture that explicitly rejects DAS in favor of algebraic dispersal proofs~\cite{Espresso:Hotshot_and_Triamisu}.
NEAR provides sharding-based DA through Nightshade with a rollup-oriented blob store~\cite{NEAR_Nightshade_Scaling_Blockchain}, and Avail combines KZG polynomial commitments with DAS-enabled light clients on a Substrate-based chain~\cite{Avail_The_DA_Blockchain}.

Published comparisons of these designs are qualitative surveys~\cite{Saif_Survey_DA_Rollups, Thibault2022Rollups}, vendor material, or measurements of a single protocol~\cite{Heimbach2025BlobFee, Lee2025BlobSharing}.
We are not aware of a study that measures throughput, cost and utilization across several DA networks with one methodology, so rollup developers choosing a DA backend have no common empirical baseline to work from.

This paper makes the following contributions:

\begin{enumerate}
    \item A taxonomy that classifies DA protocols by verification paradigm (data availability sampling, verifiable information dispersal, or validator-set recovery), encoding scheme and commitment type (Section~\ref{sec:background}).

    \item A comparison of encoding redundancy, commitment schemes (KZG, NMT and Merkle) and verification mechanisms across the six protocols, based on their specifications and documentation (Section~\ref{sec:architecture}).

    \item Measurements from 90~days of mainnet data (November 2025--February 2026) for Polkadot, Ethereum, Celestia, Espresso, NEAR and Avail, covering throughput (MiB/s), cost (\$/MiB), utilization and block dynamics under common definitions (Section~\ref{sec:results}).

    \item A comparison of trust assumptions, Byzantine fault tolerance thresholds, withholding enforcement and retention policies (Section~\ref{sec:security}).

    \item The collection scripts, raw datasets and analysis code, released under Apache~2.0 (code) and CC~BY~4.0 (data).
\end{enumerate}
